\section{\href{http://dx.doi.org/10.1088/0953-4075/39/1/009}{The effects of radiative cascades on the x-ray
diagnostic lines of Fe$^{16+}$}}

\subsection*{Supervisor: Prof Connor Ballance}

We present complete collisional-radiative modelling results for the soft x-ray
emission lines of Fe$^{16+}$ in the 15 \AA–17 \AA range. These lines have been
the subject of much controversy in the astrophysical and laboratory plasma
community. Radiative transition rates are generated from fully relativistic
atomic structure calculations. Electron-impact excitation cross sections are
determined using a fully relativistic R-matrix method employing 139 coupled
atomic levels through $n = 5$. We find that, in all cases, using a simple ratio of
the collisional rate coefficient times a radiative branching factor is not sufficient
to model the widely used diagnostic line ratios. One has to include the effects
of collisional-radiative cascades in a population model to achieve accurate line
ratios. Our line ratio results agree well with several previous calculations
and reasonably well with tokamak experimental measurements, assuming a
Maxwellian electron-energy distribution. Our modelling results for four EBIT
line ratios, assuming a narrow Gaussian electron-energy distribution, are in
generally poor agreement with all four NIST measurements but are in better
agreement with the two LLNL measurements. These results suggest the need
for an investigation of the theoretical polarization calculations that are required
to interpret the EBIT line ratio measurements.

